\subsection{Reprodução técnica do diagnóstico divergente em t3}
\label{sec:reproducao_t3}

Para investigar a hipótese de permanência de código da tarefa anterior, foram
utilizados os programas de T2 e T3 exatamente como constam do formulário
(Apêndice~\ref{ape:formulario}). A T2 contém \texttt{var total = 10 + 5;} e
\texttt{print(total;}. A T3 declara \texttt{total = 12} e \texttt{quantidade = 0},
calcula \texttt{var media = total / quantidade;} e imprime \texttt{media}.
Esses programas não recuperam o texto efetivamente inserido pelo participante,
que não foi registrado.

Cinco cenários foram executados pelo programa auxiliar
\path{rust_core/examples/reproducao_t3.rs}, com chamadas sucessivas a
\texttt{process\_source} no mesmo processo. O procedimento foi repetido sobre
a versão de 11/11/2026 e a versão final, com os mesmos
resultados relevantes, apresentados no Quadro~\ref{quad:reproducao_t3}.

\begin{quadro}[htb]
\centering
\legenda[quad:reproducao_t3]{Cenários de reprodução do diagnóstico de T3 no motor}
\footnotesize
\begin{tabular}{|P{4.4cm}|P{5.6cm}|P{2.8cm}|}
\hline
Cenário & Resultado observado nas duas revisões & Rastro \\
\hline\hline
T2 sem correção & Erro sintático: \textit{Esperado ')', mas foi encontrado ';'}. & Vazio \\
\hline
T3 substitui integralmente T2, na chamada seguinte & Erro de execução: \textit{Divisão por zero}. & Dois eventos anteriores à falha \\
\hline
T2 sem correção permanece no texto e T3 é acrescentada & Mesmo erro sintático de T2; a interpretação não é iniciada. & Vazio \\
\hline
T2 corrigida, como controle & Execução concluída, com saída 15. & Dois eventos \\
\hline
T3 substitui T2 corrigida, na chamada seguinte & Erro de execução: \textit{Divisão por zero}. & Dois eventos anteriores à falha \\
\hline
\end{tabular}
\vspace{0.2cm}
\fonte{O autor (2026).}
\end{quadro}
\FloatBarrier

A hipótese de código anterior remanescente é tecnicamente plausível quando T3 é
acrescentada à T2 malformada. Contudo, a substituição integral produz o diagnóstico
de divisão por zero, mesmo após uma compilação sintaticamente inválida no mesmo
processo. O erro sintático não persistiu entre chamadas com o código substituído. A explicação mais provável para o relato é, portanto, que o programa de T3 tenha sido acrescentado ao de T2 sem correção; como o texto inserido pelo participante não foi registrado, essa hipótese não pôde ser confirmada.
